\section{Data}
\label{sec:data}

The target dataset is the Sephora Products and Skincare Reviews collection
(Kaggle ref.\ \texttt{nadyinky/sephora-products-and-skincare-reviews}):
roughly 8{,}000 products with attributes such as category, brand, price and
ingredient lists, and on the order of one million reviews carrying a rating,
a recommendation flag, and reviewer attributes including skin type and skin
tone.

\subsection{Synthetic fixture}

The repository ships a synthetic generator
(\texttt{scripts/download\_data.py}) that reproduces this schema without
requiring Kaggle credentials, so the test suite and CI run with zero external
dependencies. Products are assigned a brand, a category, a price drawn from a
category-appropriate range, and two to four ingredient tags; one tag is
deliberately drawn from a fixed list of tags each skin type responds well to
(Table~\ref{tab:skin-affinity}), so the catalogue always contains genuine
structure for every skin type to recover, not just noise.

\begin{table}[h]
  \centering
  \caption{Synthetic ingredient-affinity clusters. A real dataset would not
  have this structure made explicit; the synthetic fixture encodes it so
  that a model with the right inductive bias (shared ingredient features)
  can be shown to recover it with a bounded sample size.}
  \label{tab:skin-affinity}
  \begin{tabular}{ll}
    \toprule
    Skin type & Preferred ingredient tags \\
    \midrule
    Dry & Hyaluronic Acid, Ceramides, Squalane, Glycerin \\
    Oily & Salicylic Acid, Niacinamide, Clay, Zinc PCA \\
    Combination & Niacinamide, Hyaluronic Acid, Green Tea \\
    Normal & Vitamin C, Peptides, Antioxidant Complex \\
    Sensitive & Centella Asiatica, Oat Extract, Ceramides, Panthenol \\
    \bottomrule
  \end{tabular}
\end{table}

Each synthetic user is assigned a skin type and a skin tone, and reviews
products with a selection probability weighted by the overlap between the
product's ingredient tags and the user's skin-type cluster, mildly tilted by
a long-tailed popularity prior. Ratings are drawn around a centre that is
higher for well-matched products and lower otherwise, and a review counts as
``recommended'' only when the rating clears the implicit-feedback threshold
(Section~\ref{sec:evaluation}). Every selection decision is keyed off
\emph{features}, never off the raw user or product ID -- this is what gives
a feature-aware model real structure to recover that an ID-only model
cannot, by construction, generalize from one identity to another.

\subsection{Implicit feedback}

A review is treated as a positive interaction when
$\text{rating} \geq 4$ or the reviewer explicitly marked the product as
recommended; everything else is dropped rather than treated as an explicit
negative, following the usual implicit-feedback convention. An iterative
$k$-core filter then removes users and products with fewer than five
surviving positives, looping because removing a product can push a user
below the floor and vice versa.
